% Sections 1.1 to 1.5, translated from sv/chapters/01/theory.tex.

\section{Number sets}
\label{c1:sec:talmangder}
Mathematics is built around different sets of numbers. Each set extends the one before it.

\begin{booktable}{@{}p{7.6em}p{4.6em}LL@{}}
  \thead{Set} & \thead{Notation} & \thead{Description} & \thead{Examples} \\
  \midrule
  Natural numbers & $\mathbb{N}$ & Positive integers and zero & $0, 1, 2, 3, \ldots$ \\
  Integers & $\mathbb{Z}$ & Natural numbers and negative numbers &
    $\ldots, -2, -1, 0, 1, 2, \ldots$ \\
  Rational\newline numbers & $\mathbb{Q}$ & Numbers that can be written as $\frac{p}{q}$, where
    $p$ and $q$ are integers and $q \neq 0$ & $\frac{1}{2},\, -\frac{3}{4},\, 0.75$ \\
  Irrational\newline numbers & -- & Numbers that \emph{cannot} be written as $\frac{p}{q}$ &
    $\pi,\, \sqrt{2},\, e$ \\
  Real numbers & $\mathbb{R}$ & All rational and irrational numbers & Every number on the number
    line \\
\end{booktable}

Electrical engineering uses all of these sets. Resistance values are often rational
($R = 4.7\,\text{k}\Omega$), while constants such as $\pi$ appear in the formulas for AC systems.

\section{The number line and absolute value}
\label{c1:sec:absolutbelopp}
Every real number can be placed on a \textbf{number line}. The number's distance from the origin
(zero) is called its \textbf{absolute value} and is written $\abs{a}$:
\[
  \abs{a} =
  \begin{cases}
    a & \text{if } a \geq 0 \\
    -a & \text{if } a < 0
  \end{cases}
\]

\begin{exempel}
\[
  \abs{5} = 5, \qquad \abs{-3} = 3, \qquad \abs{0} = 0
\]
\end{exempel}

The absolute value is never negative: $\abs{a} \geq 0$.

In electrical engineering, the absolute value is used, among other things, to give the amplitude
of an AC voltage, whether the voltage is positive or negative at the moment.

\section{Order of operations}
\label{c1:sec:rakneordning}
When an expression contains several operations, they are carried out in this order:
\begin{enumerate}
  \item \textbf{Brackets}: work out the innermost first and work outwards.
  \item \textbf{Powers and roots}: work out exponents and roots.
  \item \textbf{Multiplication and division}: from left to right.
  \item \textbf{Addition and subtraction}: from left to right.
\end{enumerate}

\begin{exempel}
\[
  3 + 2 \times 4^2 - (1 + 5) = 3 + 2 \times 16 - 6 = 3 + 32 - 6 = 29
\]
\end{exempel}

\section{Fractions}
\label{c1:sec:brak}

\subsection{Simplifying}
\label{c1:sec:forenkling}
A fraction is simplified by dividing the numerator and the denominator by their greatest common
divisor (GCD):
\[
  \frac{12}{18} = \frac{12 \div 6}{18 \div 6} = \frac{2}{3}
\]

\subsection{Addition and subtraction}
\label{c1:sec:bradd}
Fractions are added and subtracted over a \textbf{common denominator}:
\[
  \frac{a}{b} + \frac{c}{d} = \frac{ad + bc}{bd}
\]

\begin{exempel}
\[
  \frac{1}{3} + \frac{1}{4} = \frac{4}{12} + \frac{3}{12} = \frac{7}{12}
\]
\end{exempel}

\subsection{Multiplication}
\label{c1:sec:brmult}
Numerator is multiplied by numerator, denominator by denominator:
\[
  \frac{a}{b} \times \frac{c}{d} = \frac{ac}{bd}
\]

\begin{exempel}
\[
  \frac{2}{3} \times \frac{3}{5} = \frac{6}{15} = \frac{2}{5}
\]
\end{exempel}

\subsection{Division}
\label{c1:sec:brdiv}
Dividing by a fraction is the same as multiplying by the fraction turned upside down, its
\emph{reciprocal}:
\[
  \frac{a}{b} \div \frac{c}{d} = \frac{a}{b} \times \frac{d}{c} = \frac{ad}{bc}
\]

\begin{exempel}
\[
  \frac{3}{4} \div \frac{3}{8} = \frac{3}{4} \times \frac{8}{3} = \frac{24}{12} = 2
\]
\end{exempel}

\subsection{Resistors in parallel}
\label{c1:sec:parallell}
In electrical engineering, fractions arise naturally in calculating resistors in parallel. The
total resistance $R_{\text{TOT}}$ of two resistors $R_1$ and $R_2$ in parallel is given by
\[
  \frac{1}{R_{\text{TOT}}} = \frac{1}{R_1} + \frac{1}{R_2}.
\]

\begin{exempel}
With $R_1 = 6\,\Omega$ and $R_2 = 3\,\Omega$,
\[
  \frac{1}{R_{\text{TOT}}} = \frac{1}{6} + \frac{1}{3} = \frac{1}{6} + \frac{2}{6} = \frac{3}{6}
  = \frac{1}{2},
\]
so $R_{\text{TOT}} = 2\,\Omega$.
\end{exempel}

\subsection{Series and parallel connection}
\label{c1:sec:serieparallell}
When several resistors sit in the same circuit, they can be replaced by a single resistor, a
so-called \textbf{equivalent resistance}.

\textbf{A series connection} means that the resistors are connected one after the other. The
equivalent resistance $R_{\text{s}}$ is the sum of the resistances:
\[
  R_{\text{s}} = R_1 + R_2
\]

\textbf{A parallel connection} means that the resistors are connected side by side, which is
written $R_1 \parallell R_2$. Besides the formula in \secref{c1:sec:parallell}, the parallel
resistance $R_{\text{p}}$ of \emph{two} resistors can be calculated as product over sum, which is
often quicker:
\[
  R_{\text{p}} = R_1 \parallell R_2 = \frac{R_1 \times R_2}{R_1 + R_2}
\]
The parallel resistance is always smaller than the smallest of the resistors in it.

\textbf{Ohm's law} relates the voltage $U$, the resistance $R$ and the current $I$:
\[
  U = R \times I
  \qquad \Leftrightarrow \qquad
  I = \frac{U}{R}
  \qquad \Leftrightarrow \qquad
  R = \frac{U}{I}
\]

\begin{aside}
\note[Tip] If the resistance is in k$\Omega$ and the voltage in V, the current comes out directly
in mA. That saves keeping track of powers of ten.
\end{aside}

\begin{exempel}[Simplifying a circuit]
The circuit below is simplified step by step until only one resistor is left. First the parallel
connection of $R_2$ and $R_3$ is replaced by $R_{\text{p}}$, then $R_1$ is added to give the total
resistance $R_{\text{TOT}}$.

\bild[\linewidth]{lectures/L01/appendix/images/circuit_simplification.png}

With $R_1 = 2\,\text{k}\Omega$, $R_2 = 12\,\text{k}\Omega$, $R_3 = 6\,\text{k}\Omega$ and
$U = 12\,\text{V}$,
\begin{gather*}
  R_{\text{p}} = \frac{12 \times 6}{12 + 6} = \frac{72}{18} = 4\,\text{k}\Omega, \\
  R_{\text{TOT}} = R_1 + R_{\text{p}} = 2 + 4 = 6\,\text{k}\Omega, \\
  I = \frac{U}{R_{\text{TOT}}} = \frac{12}{6} = 2\,\text{mA}.
\end{gather*}
The order is decisive: the parallel connection is a calculation of its own that must be finished
before $R_1$ is added, just as brackets are worked out first in the order of operations.
\end{exempel}

\section{Percentages and proportion}
\label{c1:sec:procent}
A \textbf{percentage} is another way of writing a fraction with the denominator 100:
\[
  p\,\% = \frac{p}{100}
\]
A percentage of a value $A$ is calculated as
\[
  x\,\% \text{ of } A = \frac{x}{100} \times A.
\]

\begin{exempel}
What is $30\,\%$ of $200\,\text{V}$?
\[
  \frac{30}{100} \times 200 = 60\,\text{V}
\]
\end{exempel}

\textbf{Efficiency} is a common percentage measure in electrical engineering:
\[
  \eta = \frac{P_{\text{out}}}{P_{\text{in}}} \times 100\,\%
\]

\begin{exempel}
A transformer takes in $500\,\text{W}$ and delivers $470\,\text{W}$. The efficiency is
\[
  \eta = \frac{470}{500} \times 100\,\% = 94\,\%.
\]
\end{exempel}
